% !TeX program = lualatex
% halo:begin
% {
%   "version": 1,
%   "publish": true,
%   "course": "higher-algebra-i",
%   "section": "2.3",
%   "title": "高等代数（一）2.3：线性方程组的解",
%   "categories": ["study-notes"],
%   "tags": ["mathematics"],
%   "excerpt": "线性方程组的有解判别、基础解系与非齐次解的结构，课堂消元例题，以及含参数方程组和教材习题的完整解答。"
% }
% halo:end
% 2.3为笔记自身主题编号，非教材节号；课堂日期未确认。
% 课堂来源：IMG_2687、IMG_2688、IMG_2693--2695；续整理IMG_3176--3182（3177为3176的模糊重复）。
% 教材：北大《高等代数》第五版第三章第5、6节。
% 有解判别：印刷页92--94（PDF104--106）；解的结构：印刷页94--99（PDF106--111）。
% 使用PDF：/Users/yuchenzhang/Downloads/高等代数 第五版 北大数学系前线代小组 王萼芳 石生明 Z-Library.pdf
% 原例题为教学补充；本轮按新照片增补课堂原例及明确标出的教材练习。
% 新照片题号对应北大《高等代数》第六版：习题III第20(1)题（印刷108/PDF113）、第23题（109/114），补充题III第9题（111/116）。
% IMG_3180右页四元题仅按可确认的条件整理，不补造缺失的原题编号或系数矩阵。
% IMG_3180左页第23题特解第一分量应为-2，订正照片一处误写的-1。
\RequirePackage{iftex}
\RequireLuaTeX
\documentclass[UTF8, fontset=fandol, a4paper, 12pt]{ctexart}

\usepackage[margin=2.5cm]{geometry}
\usepackage{amsmath, amssymb, amsthm}
\usepackage{enumitem}
\usepackage{fancyhdr}
\usepackage[hidelinks]{hyperref}

\newcommand{\rank}{\operatorname{rank}}
\newcommand{\vx}{\boldsymbol{x}}
\newcommand{\vb}{\boldsymbol{\beta}}
\newcommand{\vo}{\boldsymbol{0}}
\newcommand{\va}{\boldsymbol{\alpha}}
\newcommand{\vxi}{\boldsymbol{\xi}}
\newcommand{\veta}{\boldsymbol{\eta}}
\theoremstyle{definition}
\newtheorem{definition}{定义}[section]
\newtheorem{proposition}[definition]{命题}
\newtheorem{theorem}[definition]{定理}
\newtheorem{example}[definition]{例题}
\renewcommand{\proofname}{证明}
\setlist[enumerate]{leftmargin=2em, itemsep=0.2em, topsep=0.3em}
\setlength{\headheight}{16pt}
\pagestyle{fancy}
\fancyhf{}
\fancyhead[L]{\small 高等代数（一）2.3}
\fancyhead[R]{\small 线性方程组的解}
\fancyfoot[C]{\thepage}
\renewcommand{\headrulewidth}{0.3pt}
\raggedbottom

\title{线性方程组的解}
\author{Yuchen Zhang}
\date{整理日期：2026年11月10日}
\makeatletter
\renewcommand{\maketitle}{%
  \begin{center}
    {\LARGE\bfseries \@title\par}\vspace{0.6em}
    \@author\par\vspace{0.2em}
    {\small \@date\par}
  \end{center}
}
\makeatother

\begin{document}
\maketitle

\section{线性方程组与增广矩阵}

以下设 $P$ 为数域，$m,n$ 为正整数。考虑 $P$ 上的线性方程组
\begin{equation}\label{eq:system}
 \begin{cases}
 a_{11}x_1+\cdots+a_{1n}x_n=b_1,\\
 a_{21}x_1+\cdots+a_{2n}x_n=b_2,\\
 \hspace{4em}\vdots\\
 a_{m1}x_1+\cdots+a_{mn}x_n=b_m.
 \end{cases}
\end{equation}
记系数矩阵、未知向量与常数向量分别为
\[
 A=(a_{ij})_{m\times n},\qquad
 \vx=(x_1,\ldots,x_n)^{\mathsf T},\qquad
 \vb=(b_1,\ldots,b_m)^{\mathsf T}.
\]
则方程组写为 $A\vx=\vb$。记 $A=(\va_1,\ldots,\va_n)$，又可写成
\[
 x_1\va_1+\cdots+x_n\va_n=\vb.
\]
因此，方程组有解，当且仅当 $\vb$ 可由 $A$ 的列向量组线性表示。
这里 $\va_j,\vb\in P^m$，而解向量 $\vx\in P^n$。

\begin{definition}[增广矩阵]
在系数矩阵 $A$ 的右侧添加常数列 $\vb$，得到的 $m\times(n+1)$ 矩阵
\[
 \overline A=(A\mid\vb)
 =\left(\begin{array}{ccc|c}
 a_{11}&\cdots&a_{1n}&b_1\\
 \vdots&&\vdots&\vdots\\
 a_{m1}&\cdots&a_{mn}&b_m
 \end{array}\right)
\]
称为方程组\eqref{eq:system}的增广矩阵。
\end{definition}

对增广矩阵作初等行变换，相当于对整个方程作同解变换，常数项必须同时变换。
交换两行、某行乘非零数、将一行的倍数加到另一行，这三种操作都可逆，因此保持解集不变。
初等列变换虽然保持矩阵的秩，却一般不保持原未知量下的解；交换未知量对应的列时，须同时记录变量次序。

\section{线性方程组有解的判别}

\begin{theorem}[有解判别定理]
线性方程组 $A\vx=\vb$ 有解，当且仅当
\[
 \boxed{\rank A=\rank(A\mid\vb).}
\]
\end{theorem}
\begin{proof}
先证必要性。若有解，则 $\vb=\sum_{j=1}^n x_j\va_j$，故向量组
$\va_1,\ldots,\va_n$ 与 $\va_1,\ldots,\va_n,\vb$ 等价，两组的秩相同。
由矩阵的秩等于列秩，得到 $\rank A=\rank(A\mid\vb)$。

再证充分性。设两矩阵的秩均为 $r$。若 $r=0$，则 $A=0$ 且 $\vb=0$，任意 $\vx\in P^n$ 都是解。
以下设 $r>0$，取 $A$ 的列向量组的一个极大无关组
$\va_{i_1},\ldots,\va_{i_r}$。
若 $\vb$ 不能由它们线性表示，则
$\va_{i_1},\ldots,\va_{i_r},\vb$ 线性无关，
从而 $\rank(A\mid\vb)\ge r+1$，矛盾。
故 $\vb$ 可由 $A$ 的列向量组表示，表示系数给出方程组的解。
\end{proof}

在 $A$ 中只增加一列，秩至多增加 $1$，所以
\[
 \rank(A\mid\vb)=\rank A\quad\text{或}\quad\rank A+1.
\]
若对增广矩阵行消元后出现
\[
 (0\ \cdots\ 0\mid c),\qquad c\ne0,
\]
则对应方程为 $0=c$，原方程组无解。此时常数列产生一个额外主元，增广秩比系数秩大 $1$。
而行 $(0\ \cdots\ 0\mid0)$ 只对应恒等式 $0=0$，不增加约束。

\subsection{主元未知量与自由未知量}

\begin{definition}
将增广矩阵用初等行变换化为行阶梯形。在方程组有解时，系数部分主元所在列对应的未知量
称为\textbf{主元未知量}，其余未知量称为\textbf{自由未知量}。
\end{definition}
主元列是非零行的首个非零元素所在的列，编号不必连续。
自由未知量可以任意指定；指定以后，由末行向上回代，主元未知量随之唯一确定。

\begin{theorem}[解的个数]
设 $A$ 有 $n$ 列，$r=\rank A$，则：
\begin{enumerate}
\item 若 $\rank A<\rank\overline A$，方程组无解；
\item 若 $\rank A=\rank\overline A=n$，方程组有唯一解；
\item 若 $\rank A=\rank\overline A=r<n$，方程组有无穷多解，且自由未知量有 $n-r$ 个。
\end{enumerate}
\end{theorem}
\begin{proof}
第一项即有解判别定理。以下设方程组有解。初等行变换保持秩，
阶梯形系数矩阵有 $r$ 个非零行，也就有 $r$ 个主元。
因此有 $r$ 个主元未知量、$n-r$ 个自由未知量。

若 $r=n$，没有自由未知量，逐行回代得到唯一解。
若 $r<n$，任意给定自由未知量的值，都能回代得到唯一一组主元分量。
不同的自由分量给出不同的解，而数域 $P$ 含无穷多个数，故解有无穷多个。
\end{proof}
这里与秩比较的是未知量个数 $n$，不是方程个数 $m$。
仅有 $\rank A=n$ 只能保证\textbf{有解时解唯一}，不能保证非齐次方程组一定有解。
例如 $A=(1,0)^{\mathsf T}$、$\vb=(0,1)^{\mathsf T}$ 时，
方程组为 $x_1=0,\ 0=1$，虽 $\rank A=n=1$，仍无解。

\begin{example}\label{ex:nonhom}
解线性方程组
\[
 \begin{cases}
 x_1+2x_2-x_3+x_4=1,\\
 2x_1+4x_2+x_3-x_4=5,\\
 x_1+2x_2+2x_3-2x_4=4.
 \end{cases}
\]
\end{example}
\begin{proof}[解]
对增广矩阵作初等行变换：
\begin{align*}
\left(\begin{array}{rrrr|r}1&2&-1&1&1\\2&4&1&-1&5\\1&2&2&-2&4\end{array}\right)
&\xrightarrow[ R_3\leftarrow R_3-R_1]{R_2\leftarrow R_2-2R_1}
\left(\begin{array}{rrrr|r}1&2&-1&1&1\\0&0&3&-3&3\\0&0&3&-3&3\end{array}\right)\\
&\xrightarrow{R_3\leftarrow R_3-R_2}
\left(\begin{array}{rrrr|r}1&2&-1&1&1\\0&0&3&-3&3\\0&0&0&0&0\end{array}\right)\\
&\xrightarrow{R_2\leftarrow R_2/3}
\left(\begin{array}{rrrr|r}1&2&-1&1&1\\0&0&1&-1&1\\0&0&0&0&0\end{array}\right)\\
&\xrightarrow{R_1\leftarrow R_1+R_2}
\left(\begin{array}{rrrr|r}1&2&0&0&2\\0&0&1&-1&1\\0&0&0&0&0\end{array}\right).
\end{align*}
故 $\rank A=\rank\overline A=2<4$，方程组有无穷多解。
主元位于第 $1,3$ 列，故 $x_2,x_4$ 为自由未知量。令 $x_2=s,x_4=t$，由
\[
 x_1+2x_2=2,\qquad x_3-x_4=1
\]
得到全部解
\[
 \boxed{\vx=(2-2s,s,1+t,t)^{\mathsf T},\qquad s,t\in P.}
\]
任意 $s,t$ 都给出一个解；反之，任意解的 $x_2,x_4$ 一旦给定，$x_1,x_3$ 即由上式确定，故此表示没有遗漏。
\end{proof}

\newpage
\section{齐次线性方程组与基础解系}

\begin{definition}[齐次线性方程组]
常数项全为零的方程组 $A\vx=\vo$ 称为齐次线性方程组。
向量 $\vx=\vo$ 总是它的解，称为零解；分量不全为零的解称为非零解。
\end{definition}

\begin{theorem}
有 $n$ 个未知量的齐次线性方程组 $A\vx=\vo$ 有非零解，当且仅当 $\rank A<n$。
\end{theorem}
\begin{proof}
添加零列不改变秩，故齐次方程组总满足有解条件。
若 $\rank A=n$，由解的个数定理，零解是唯一解。
若 $\rank A<n$，至少存在一个自由未知量。将其中一个取为 $1$，其余自由未知量取为 $0$，
再回代求主元未知量，得到的解含有分量 $1$，因而是非零解。
\end{proof}

\begin{proposition}
若方程个数 $m$ 少于未知量个数 $n$，则齐次方程组必有非零解。
若 $A$ 为 $n$ 阶方阵，则齐次方程组有非零解当且仅当 $\det A=0$。
\end{proposition}
\begin{proof}
当 $m<n$ 时，$\rank A\le m<n$，由上一定理得到第一项。
当 $A$ 为 $n$ 阶方阵时，$\rank A<n$ 当且仅当 $\det A=0$，得到第二项。
\end{proof}

\begin{proposition}[齐次解的线性组合]
若 $A\boldsymbol u=\vo$、$A\boldsymbol v=\vo$，则对任意 $s,t\in P$，
$s\boldsymbol u+t\boldsymbol v$ 仍为齐次方程组的解。
\end{proposition}
\begin{proof}
由矩阵乘向量的线性性，
\[
 A(s\boldsymbol u+t\boldsymbol v)=sA\boldsymbol u+tA\boldsymbol v=\vo.
\]
\end{proof}
有限个齐次解的线性组合也仍是解。因此，可以选取适当的一组解，用它们的线性组合表示全部解。

\begin{definition}[基础解系]
设齐次方程组有非零解。若它的一组解 $\vxi_1,\ldots,\vxi_k$ 线性无关，
且方程组的每个解都能由它们线性表示，则称这组解为方程组的\textbf{基础解系}。
此时全部解为
\[
 \vx=c_1\vxi_1+\cdots+c_k\vxi_k,\qquad c_1,\ldots,c_k\in P.
\]
\end{definition}
基础解系的各向量必须同时满足：是原齐次方程组的解、线性无关、能表示全部解。
沿用这里的定义，只有零解时不列基础解系。

\subsection{基础解系的存在与构造}

\begin{theorem}
设 $A$ 有 $n$ 列，$\rank A=r<n$。则 $A\vx=\vo$ 存在基础解系，
且任意基础解系都恰含 $n-r$ 个解向量。
\end{theorem}
\begin{proof}
将 $A$ 化为行阶梯形，设自由未知量为 $x_{j_1},\ldots,x_{j_k}$，其中 $k=n-r$。
分别令这 $k$ 个自由分量取值为
\[
 (1,0,\ldots,0),\quad(0,1,\ldots,0),\quad\ldots,\quad(0,\ldots,0,1),
\]
并回代求其余分量，得到 $k$ 个完整的 $n$ 维解向量 $\vxi_1,\ldots,\vxi_k$。
下面证明它们构成基础解系。

\textbf{第一步：线性无关。}设 $\sum_{i=1}^k c_i\vxi_i=\vo$。
对每个 $\ell=1,\ldots,k$，取第 $j_\ell$ 个分量，左边为 $c_\ell$，右边为 $0$，
故 $c_1=\cdots=c_k=0$。

\textbf{第二步：表示全部解。}任取齐次解 $\vx$，设其自由分量为
$x_{j_1}=t_1,\ldots,x_{j_k}=t_k$。向量
$\boldsymbol y=\sum_{i=1}^k t_i\vxi_i$ 也是齐次解，并且与 $\vx$ 具有相同的自由分量。
自由分量确定后，主元分量由回代唯一确定，故 $\vx=\boldsymbol y$。
因此每个齐次解都能由 $\vxi_1,\ldots,\vxi_k$ 表出。

最后，任意另一基础解系与此组都能表示全部解，所以二者等价；又都线性无关，
由等价无关组的向量个数相同，另一基础解系也含 $k=n-r$ 个向量。
\end{proof}
自由未知量取单位值时，只有\textbf{自由分量}按单位向量赋值，主元分量仍须回代计算。
每个基础解向量有 $n$ 个分量，而基础解系共有 $n-r$ 个向量。
当 $r=0$ 时，$A=0$，可取 $P^n$ 的标准单位向量组为基础解系。

\begin{example}\label{ex:hom}
求齐次方程组
\[
 \begin{cases}
 x_1+2x_2-x_3+x_4=0,\\
 2x_1+4x_2+x_3-x_4=0,\\
 x_1+2x_2+2x_3-2x_4=0
 \end{cases}
\]
的一个基础解系及全部解。
\end{example}
\begin{proof}[解]
系数矩阵与例\ref{ex:nonhom}相同，作相同的初等行变换，得到
\[
 \begin{pmatrix}1&2&0&0\\0&0&1&-1\\0&0&0&0\end{pmatrix}.
\]
故 $x_1=-2x_2$、$x_3=x_4$，自由未知量为 $x_2,x_4$。
分别取 $(x_2,x_4)=(1,0)$、$(0,1)$，得到
\[
 \vxi_1=(-2,1,0,0)^{\mathsf T},\qquad
 \vxi_2=(0,0,1,1)^{\mathsf T}.
\]
两向量均满足齐次方程组。若 $c_1\vxi_1+c_2\vxi_2=\vo$，
考察第 $2,4$ 个分量得 $c_1=c_2=0$，所以二者无关。
任意解可令 $x_2=s,x_4=t$，写成
\[
 \vx=(-2s,s,t,t)^{\mathsf T}=s\vxi_1+t\vxi_2.
\]
故 $\vxi_1,\vxi_2$ 构成一个基础解系，全部解为
\[
 \boxed{\vx=s(-2,1,0,0)^{\mathsf T}+t(0,0,1,1)^{\mathsf T},\qquad s,t\in P.}
\]
\end{proof}
基础解系与系数矩阵的列极大无关组是不同对象：前者由 $n-r$ 个 $n$ 维解向量组成，
后者在 $r>0$ 时由原矩阵中 $r$ 个 $m$ 维列向量组成，不能互相替代。

\subsection{课堂例题：五元齐次方程组}

下面按自由分量取单位值的方法，求课堂五元方程组的基础解系。

\begin{example}\label{ex:class-hom}
求齐次方程组的一个基础解系及全部解：
\[
 \begin{cases}
 x_1+x_2+x_3+x_4+2x_5=0,\\
 x_1+x_2+2x_3+2x_4+3x_5=0,\\
 2x_1+2x_2+2x_3+3x_4+5x_5=0,\\
 4x_1+4x_2+4x_3+5x_4+9x_5=0.
 \end{cases}
\]
\end{example}
\begin{proof}[解]
第二、三、四行先分别减去第一行的 $1,2,4$ 倍，第四行再减去第三行，得到
\begin{align*}
 A&=\begin{pmatrix}1&1&1&1&2\\1&1&2&2&3\\2&2&2&3&5\\4&4&4&5&9\end{pmatrix}\\
 &\longrightarrow
 \begin{pmatrix}1&1&1&1&2\\0&0&1&1&1\\0&0&0&1&1\\0&0&0&1&1\end{pmatrix}
 \longrightarrow
 \begin{pmatrix}1&1&1&1&2\\0&0&1&1&1\\0&0&0&1&1\\0&0&0&0&0\end{pmatrix}.
\end{align*}
主元位于第 $1,3,4$ 列，故 $\rank A=3$，自由未知量为 $x_2,x_5$。
令 $x_2=s,x_5=t$，由末行向上回代，依次得到
\[
 x_4=-t,\qquad x_3=0,\qquad x_1=-s-t.
\]
所以全部解为
\[
 \vx=s\underbrace{(-1,1,0,0,0)^{\mathsf T}}_{\vxi_1}
       +t\underbrace{(-1,0,0,-1,1)^{\mathsf T}}_{\vxi_2},\qquad s,t\in P.
\]
两个向量在第 $2,5$ 个分量上的取值分别为 $(1,0)$、$(0,1)$，所以线性无关；
上式又表示了全部解，故 $\vxi_1,\vxi_2$ 构成一个基础解系。
\end{proof}

\section{非齐次线性方程组解的结构}

设 $\vb\ne\vo$。把非齐次方程组 $A\vx=\vb$ 的常数项全换为零，
得到的齐次方程组 $A\vx=\vo$ 称为它的\textbf{导出组}。

\begin{proposition}
非齐次方程组的任意两个解之差是导出组的解；
非齐次方程组的一个解与导出组的一个解之和，仍是非齐次方程组的解。
\end{proposition}
\begin{proof}
若 $A\veta_1=A\veta_2=\vb$，则
\[
 A(\veta_1-\veta_2)=\vb-\vb=\vo.
\]
若 $A\veta=\vb$、$A\vxi=\vo$，则
\[
 A(\veta+\vxi)=\vb+\vo=\vb.
\]
\end{proof}

\begin{theorem}[非齐次方程组的通解]
设 $A\vx=\vb$ 有解，$\veta$ 是它的一个特解，则全部解为
\[
 \boxed{\vx=\veta+\vxi,\qquad \vxi\text{ 取遍导出组的全部解}.}
\]
若 $\rank A=r<n$，且 $\vxi_1,\ldots,\vxi_{n-r}$ 是导出组的一个基础解系，则
\[
 \boxed{\vx=\veta+c_1\vxi_1+\cdots+c_{n-r}\vxi_{n-r},\qquad c_i\in P.}
\]
若 $r=n$，则原方程组只有唯一解 $\veta$。
\end{theorem}
\begin{proof}
任取原方程组的解 $\vx$，有
\[
 \vx=\veta+(\vx-\veta).
\]
由上一命题，$\vx-\veta$ 是导出组的解，故原方程组的每个解都具有所述形式。
反过来，对任意导出组的解 $\vxi$，$\veta+\vxi$ 都是原方程组的解。
所以此表示恰好给出全部解。
当 $r<n$ 时，用基础解系表示 $\vxi$ 得到第二式；当 $r=n$ 时导出组只有零解，故 $\vx=\veta$。
\end{proof}
特解与基础解系都不必唯一。更换特解时，新旧特解之差属于导出组的解集，因此通解表示的集合不变。
对非齐次方程组，任意两个解之和一般不是解，因为其右端变为 $2\vb$；
这与齐次解对线性组合封闭的性质不同。

\begin{example}
用一个特解与导出组的基础解系表示例\ref{ex:nonhom}中方程组的全部解。
\end{example}
\begin{proof}[解]
由例\ref{ex:nonhom}，取自由未知量 $x_2=x_4=0$，得到一个特解
$\veta=(2,0,1,0)^{\mathsf T}$。
由例\ref{ex:hom}，导出组的基础解系可取
$\vxi_1=(-2,1,0,0)^{\mathsf T}$、$\vxi_2=(0,0,1,1)^{\mathsf T}$。
故全部解为
\[
 \boxed{\vx=
 \begin{pmatrix}2\\0\\1\\0\end{pmatrix}
 +s\begin{pmatrix}-2\\1\\0\\0\end{pmatrix}
 +t\begin{pmatrix}0\\0\\1\\1\end{pmatrix},\qquad s,t\in P.}
\]
代入原系数矩阵，有
\[
 A\veta=(1,5,4)^{\mathsf T}=\vb,\qquad A\vxi_1=A\vxi_2=\vo,
\]
与解的结构定理一致。
\end{proof}

\subsection{课堂例题：同一系数矩阵的非齐次方程组}

保留例\ref{ex:class-hom}的系数矩阵，只把常数列改为 $(1,2,1,3)^{\mathsf T}$。

\begin{example}
求方程组的全部解：
\[
 \begin{cases}
 x_1+x_2+x_3+x_4+2x_5=1,\\
 x_1+x_2+2x_3+2x_4+3x_5=2,\\
 2x_1+2x_2+2x_3+3x_4+5x_5=1,\\
 4x_1+4x_2+4x_3+5x_4+9x_5=3.
 \end{cases}
\]
\end{example}
\begin{proof}[解]
对增广矩阵作与例\ref{ex:class-hom}相同的行变换，再消去主元上方的元素：
\begin{align*}
 \overline A
 &=\left(\begin{array}{rrrrr|r}1&1&1&1&2&1\\1&1&2&2&3&2\\2&2&2&3&5&1\\4&4&4&5&9&3\end{array}\right)\\
 &\longrightarrow
 \left(\begin{array}{rrrrr|r}1&1&1&1&2&1\\0&0&1&1&1&1\\0&0&0&1&1&-1\\0&0&0&0&0&0\end{array}\right)\\
 &\longrightarrow
 \left(\begin{array}{rrrrr|r}1&1&0&0&1&0\\0&0&1&0&0&2\\0&0&0&1&1&-1\\0&0&0&0&0&0\end{array}\right).
\end{align*}
系数秩与增广秩均为 $3$。令自由未知量 $x_2=x_5=0$，得到一个特解
\[
 \veta=(0,0,2,-1,0)^{\mathsf T}.
\]
导出组正是例\ref{ex:class-hom}的齐次方程组，因此全部解为
\[
 \boxed{\vx=(0,0,2,-1,0)^{\mathsf T}
 +s(-1,1,0,0,0)^{\mathsf T}
 +t(-1,0,0,-1,1)^{\mathsf T},\quad s,t\in P.}
\]
\end{proof}

\subsection{课堂例题：三式四元方程组}

\begin{example}
求方程组的全部解：
\[
 \begin{cases}
 9x_1-3x_2+5x_3+6x_4=4,\\
 6x_1-2x_2+3x_3+x_4=5,\\
 3x_1-x_2+3x_3+14x_4=-8.
 \end{cases}
\]
\end{example}
\begin{proof}[解]
交换第一、三行，再消去第一列，得到
\begin{align*}
 \left(\begin{array}{rrrr|r}9&-3&5&6&4\\6&-2&3&1&5\\3&-1&3&14&-8\end{array}\right)
 &\longrightarrow
 \left(\begin{array}{rrrr|r}3&-1&3&14&-8\\0&0&-3&-27&21\\0&0&-4&-36&28\end{array}\right)\\
 &\longrightarrow
 \left(\begin{array}{rrrr|r}3&-1&3&14&-8\\0&0&1&9&-7\\0&0&0&0&0\end{array}\right).
\end{align*}
两个主元在第 $1,3$ 列，自由未知量为 $x_2,x_4$。
令 $x_2=s,x_4=t$，回代得到全部解：
\[
 \boxed{\vx=
 \begin{pmatrix}13/3\\0\\-7\\0\end{pmatrix}
 +s\begin{pmatrix}1/3\\1\\0\\0\end{pmatrix}
 +t\begin{pmatrix}13/3\\0\\-9\\1\end{pmatrix},\qquad s,t\in P.}
\]
取 $s=t=0$ 可得一个特解；两个齐次向量在第 $2,4$ 个分量上的取值分别为 $(1,0)$、$(0,1)$，
由前述构造方法，它们组成导出组的一个基础解系。
\end{proof}

\subsection{课堂例题：由已知解求全部解}

\begin{example}
设某个四元非齐次方程组 $A\vx=\vb$ 的系数矩阵满足 $\rank A=3$。
已知 $\veta_1,\veta_2,\veta_3$ 都是它的解，且
\[
 \veta_1=(1,-2,-3,4)^{\mathsf T},\qquad
 5\veta_2-3\veta_3=(2,0,0,8)^{\mathsf T}.
\]
求这个方程组的全部解。
\end{example}
\begin{proof}[解]
由 $A\veta_2=A\veta_3=\vb$，有
\[
 A(5\veta_2-3\veta_3)=2\vb.
\]
所以 $\veta_0=\tfrac12(5\veta_2-3\veta_3)=(1,0,0,4)^{\mathsf T}$ 也是一个特解。
两个特解之差
\[
 \vxi=\veta_0-\veta_1=(0,2,3,0)^{\mathsf T}
\]
是导出组的非零解。又因 $n-\rank A=4-3=1$，导出组的基础解系只含一个向量，
故 $\vxi$ 本身组成基础解系，原方程组的全部解为
\[
 \boxed{\vx=(1,-2,-3,4)^{\mathsf T}+c(0,2,3,0)^{\mathsf T},\qquad c\in P.}
\]
这里不能直接把 $5\veta_2-3\veta_3$ 当作原方程组的解：它的系数和为 $2$，须先除以 $2$。
\end{proof}

\section{含参数的线性方程组}

\begin{example}
讨论实参数 $a$ 取不同值时，下列方程组的解：
\[
 \begin{cases}
 x+y+z=1,\\
 x+ay+z=a,\\
 x+y+a^2z=a.
 \end{cases}
\]
\end{example}
\begin{proof}[解]
第二、三式分别减去第一式，得到同解方程组
\[
 x+y+z=1,\qquad (a-1)y=a-1,\qquad (a^2-1)z=a-1.
\]
由于 $a-1$、$a^2-1$ 可能为零，分以下三种情形。
\begin{enumerate}
\item 若 $a\ne\pm1$，则 $y=1$、$z=1/(a+1)$，回代得 $x=-1/(a+1)$。
此时系数秩与增广秩均为 $3$，唯一解为
\[
 (x,y,z)=\left(-\frac1{a+1},1,\frac1{a+1}\right).
\]
\item 若 $a=1$，后两式均为 $0=0$，只剩 $x+y+z=1$。
系数秩与增广秩均为 $1$，令 $y=s,z=t$，得全部解
\[
 (x,y,z)^{\mathsf T}=(1,0,0)^{\mathsf T}
 +s(-1,1,0)^{\mathsf T}+t(-1,0,1)^{\mathsf T},\qquad s,t\in\mathbb R.
\]
\item 若 $a=-1$，第三式变成 $0=-2$，矛盾，故无解。
此时系数矩阵秩为 $2$，增广矩阵秩为 $3$。
\end{enumerate}
\end{proof}
该方程组的系数行列式为 $(a-1)^2(a+1)$。
行列式非零时解唯一；行列式为零时，$a=1$ 给出无穷多解，而 $a=-1$ 给出无解。
因此，对非齐次方程组，仅由系数行列式为零不能判定是否有解，还须比较系数秩与增广秩。

\section{课后作业}

本节题号按北大《高等代数》第六版标注。

\subsection{习题 III 第20题（1）：对角元含参数的方程组}

讨论 $\lambda\in P$ 取什么值时下列方程组有解，并求解：
\[
 \begin{cases}
 \lambda x_1+x_2+x_3=1,\\
 x_1+\lambda x_2+x_3=\lambda,\\
 x_1+x_2+\lambda x_3=\lambda^2.
 \end{cases}
\]
\begin{proof}[解]
先找出系数行列式为零的参数值。沿第一行展开并因式分解，得到
\[
 D=\begin{vmatrix}\lambda&1&1\\1&\lambda&1\\1&1&\lambda\end{vmatrix}
 =\lambda^3-3\lambda+2=(\lambda-1)^2(\lambda+2).
\]
因此分三种情形。
\begin{enumerate}
 \item 若 $\lambda\ne1,-2$，则 $D\ne0$，由克拉默法则，方程组有唯一解。
 将第一、二、三列分别换成常数列，计算得
 \begin{align*}
 D_1&=-(\lambda-1)^2(\lambda+1),\\
 D_2&=(\lambda-1)^2,\\
 D_3&=(\lambda-1)^2(\lambda+1)^2.
 \end{align*}
 所以唯一解为
 \[
  \boxed{(x_1,x_2,x_3)=
  \left(-\frac{\lambda+1}{\lambda+2},\frac1{\lambda+2},
  \frac{(\lambda+1)^2}{\lambda+2}\right).}
 \]
 \item 若 $\lambda=1$，三式都化为 $x_1+x_2+x_3=1$。
 令 $x_2=s,x_3=t$，得到全部解
 \[
  \vx=(1,0,0)^{\mathsf T}+s(-1,1,0)^{\mathsf T}
  +t(-1,0,1)^{\mathsf T},\qquad s,t\in P.
 \]
 \item 若 $\lambda=-2$，将三式相加，左端为 $0$，右端为
 $1-2+4=3$，矛盾，故无解。
\end{enumerate}
两个使行列式为零的参数值，分别给出无穷多解与无解，不能省去这两种情形的检查。
\end{proof}

\subsection{习题 III 第23题：先讨论有解条件，再求通解}

$a,b\in P$ 取什么值时，线性方程组
\[
 \begin{cases}
 x_1+x_2+x_3+x_4+x_5=1,\\
 3x_1+2x_2+x_3+x_4-3x_5=a,\\
 x_2+2x_3+2x_4+6x_5=3,\\
 5x_1+4x_2+3x_3+3x_4-x_5=b
 \end{cases}
\]
有解？在有解的情形下，求它的全部解。
\newpage
\begin{proof}[解]
先消去第一列，再用第三行消去第二、四行的非零系数，最后将第三行换到第二行，得到
\begin{align*}
 \overline A
 &=\left(\begin{array}{rrrrr|r}1&1&1&1&1&1\\3&2&1&1&-3&a\\0&1&2&2&6&3\\5&4&3&3&-1&b\end{array}\right)\\
 &\longrightarrow
 \left(\begin{array}{rrrrr|r}1&1&1&1&1&1\\0&-1&-2&-2&-6&a-3\\0&1&2&2&6&3\\0&-1&-2&-2&-6&b-5\end{array}\right)\\
 &\longrightarrow
 \left(\begin{array}{rrrrr|r}1&1&1&1&1&1\\0&1&2&2&6&3\\0&0&0&0&0&a\\0&0&0&0&0&b-2\end{array}\right).
\end{align*}
末两行分别对应 $0=a$、$0=b-2$，故有解的充要条件为
\[
 \boxed{a=0,\qquad b=2.}
\]
此时系数秩与增广秩均为 $2$，有三个自由未知量 $x_3,x_4,x_5$。
令 $x_3=s,x_4=t,x_5=u$，由第二行及第一行依次得到
\[
 x_2=3-2s-2t-6u,\qquad x_1=-2+s+t+5u.
\]
所以全部解为
\begin{align*}
 \vx={}&(-2,3,0,0,0)^{\mathsf T}
 +s(1,-2,1,0,0)^{\mathsf T}\\
 &+t(1,-2,0,1,0)^{\mathsf T}
 +u(5,-6,0,0,1)^{\mathsf T},\qquad s,t,u\in P.
\end{align*}
其中 $(-2,3,0,0,0)^{\mathsf T}$ 是一个特解；后三个向量均为导出组的解，
且它们在第 $3,4,5$ 个分量上的取值组成标准单位向量组，因此构成导出组的一个基础解系。
\end{proof}

\subsection{补充题 III 第9题：系数和为1的解表示}

设 $\veta_0$ 是线性方程组 $A\vx=\vb$ 的一个解，
$\veta_1,\ldots,\veta_t$ 是它的导出组的一个基础解系。令
\[
 \boldsymbol\gamma_1=\veta_0,\qquad
 \boldsymbol\gamma_{i+1}=\veta_i+\veta_0\quad(i=1,\ldots,t).
\]
证明：原方程组的任意一个解 $\boldsymbol\gamma$ 都可以表示成
\[
 \boldsymbol\gamma=u_1\boldsymbol\gamma_1+\cdots+u_{t+1}\boldsymbol\gamma_{t+1},
 \qquad u_1+\cdots+u_{t+1}=1.
\]
\begin{proof}
任取原方程组的解 $\boldsymbol\gamma$。由非齐次解的结构定理，存在 $c_1,\ldots,c_t\in P$，使
\[
 \boldsymbol\gamma=\veta_0+\sum_{i=1}^{t}c_i\veta_i.
\]
将 $\veta_i=\boldsymbol\gamma_{i+1}-\boldsymbol\gamma_1$ 代入，得
\begin{align*}
 \boldsymbol\gamma
 &=\boldsymbol\gamma_1+
 \sum_{i=1}^{t}c_i(\boldsymbol\gamma_{i+1}-\boldsymbol\gamma_1)\\
 &=\left(1-\sum_{i=1}^{t}c_i\right)\boldsymbol\gamma_1
 +\sum_{i=1}^{t}c_i\boldsymbol\gamma_{i+1}.
\end{align*}
取 $u_1=1-\sum_{i=1}^{t}c_i$、$u_{i+1}=c_i$，便得到所需表示，且这些系数之和为 $1$。
\end{proof}
反过来，因为每个 $\boldsymbol\gamma_i$ 都满足 $A\boldsymbol\gamma_i=\vb$，
任何系数和为 $1$ 的线性组合也满足
\[
 A\left(\sum_{i=1}^{t+1}u_i\boldsymbol\gamma_i\right)
 =\left(\sum_{i=1}^{t+1}u_i\right)\vb=\vb.
\]
因此，这些系数和为 $1$ 的组合恰好给出原方程组的全部解。
若导出组只有零解，原方程组的解唯一，直接取 $u_1=1$ 即可。

\end{document}
